% Introduction
% Covers: problem statement, motivation, research gap, RQs, contributions, paper structure

\label{sec:introduction}

\subsection{Problem Statement}

Project-based courses across disciplines share a common administrative challenge:
the course lifecycle unfolds across distinct organizational phases -- application, team formation, iterative project work, and formative assessment -- yet no existing course administration tool provides a structured model for managing these phases as first-class entities.
Instructors of project-based courses currently rely on fragile combinations of spreadsheets, email, and general-purpose \acrfull{lms} plugins, resulting in recurring manual effort, limited process consistency, and poor lifecycle transparency for students and instructors alike.

\subsection{Motivation}

The absence of purpose-built administration tooling has concrete consequences:
teams are formed through manual spreadsheet coordination, assessment criteria are communicated informally rather than enforced through tooling, and each course iteration requires instructors to reconstruct the same administrative workflows from scratch.
These limitations motivate a platform that models the course lifecycle explicitly, making phases, transitions, and evaluation criteria first-class design elements.

\subsection{Research Gap}

Existing \acrshort{lms} platforms such as Moodle~\cite{moodle} and Canvas model courses as flat repositories of content and assignments, providing no native support for phase-specific workflows or lifecycle-driven administration.
Specialized computing education platforms such as Artemis~\cite{artemis} and GitHub Classroom~\cite{github-classroom} optimize for individual exercise submission and automated grading, addressing a different administrative problem.
No existing open-source platform provides a composable, phase-based abstraction for managing the full administrative lifecycle of project-based courses.

\subsection{Research Questions}

This paper addresses the following research questions:
\begin{enumerate}
    \item \textbf{RQ1:} Does the phase-based course meta-model in PROMPT adequately capture the administrative structure of project-based courses?
    \item \textbf{RQ2:} How do instructors and \acrshortpl{ta} perceive the utility and usability of PROMPT for managing project-based course administration?
    \item \textbf{RQ3:} How effectively does the plugin architecture support institution-specific and discipline-specific extensions without modifying the platform core?
\end{enumerate}

\subsection{Contributions}

The main contributions of this paper are:
\begin{itemize}
    \item A phase-based course meta-model that treats the course lifecycle as an ordered sequence of composable, independently operable phase instances -- a reusable abstraction for any project-based course regardless of discipline.
    \item PROMPT, an open-source course administration platform implementing the meta-model, with four production-ready default phase implementations (Application, Interview, Team Allocation, and Assessment) and a plugin architecture for custom extensions.
    \item An empirical evaluation via structured instructor interviews and a student \acrshort{sus} survey, complemented by a feature comparison with existing tools and multi-semester deployment evidence.
\end{itemize}

\subsection{Paper Structure}

\autoref{sec:background} reviews existing course administration tools and establishes the research gap.
\autoref{sec:prompt} describes the PROMPT platform, its phase-based meta-model, default phase implementations, and system architecture.
\autoref{sec:evaluation} presents the evaluation methodology and results addressing the three research questions.
\autoref{sec:discussion} interprets the findings, discusses broader applicability, and acknowledges limitations.
\autoref{sec:conclusion} summarizes the contributions and outlines future work.
